\subsection{Results}
\label{sec:exp1}

We trained PPO and A2C across all combinations of three design choices: the
reward function (the \emph{base} task reward vs.\ the \emph{tuned} shaped
reward), the observation encoding (raw $[\theta,\omega]$ vs.\ the
wrap-invariant $[\sin\theta,\cos\theta,\omega]$), and the disturbance setting
(nominal vs.\ \emph{robust} training with actuator noise and dropped commands).
Each configuration was run with \textbf{six random seeds} (42--47); every run
was capped at $10^6$ steps with early stopping, which also halted training as
soon as the evaluation return reached $180$ (i.e.\ $0.9\times$ the maximum
attainable return of $200$). We report, aggregated as mean~$\pm$~std over seeds,
the environment steps trained, the wall-clock training time, and the mean return
over 10 evaluation episodes (Table~\ref{tab:exp1}); the design is paired
(matched seeds), so the comparisons in Table~\ref{tab:exp1_tests} are run within
seed blocks using the Wilcoxon signed-rank test with bootstrap confidence
intervals.

\begin{table}[t]
\centering
\caption{Swing-up results over 6 seeds (42--47) per configuration. ``Return'' is mean $\pm$ std over seeds of the mean return over 10 eval episodes. ``$n_{180}$'' is the number of seeds that reached the early-stopping cap of 180; \textbf{bold} return marks censored configurations (return is then a floor, not a ceiling).}
\label{tab:exp1}
\begin{tabular}{llll rr r c}
\toprule
Algo & Reward & Obs & Dist. & Steps & Time [min] & Return & $n_{180}$ \\
\midrule
PPO & tuned & sincos & nom. & 167k & 1.4 & \textbf{181.1 $\pm$ 0.7} & 6/6 \\
PPO & base & sincos & rob. & 233k & 2.0 & \textbf{181.1 $\pm$ 0.3} & 6/6 \\
PPO & base & sincos & nom. & 192k & 1.8 & \textbf{180.8 $\pm$ 0.4} & 6/6 \\
PPO & tuned & sincos & rob. & 342k & 3.2 & \textbf{178.3 $\pm$ 4.5} & 4/6 \\
PPO & base & raw & nom. & 496k & 5.4 & 163.4 $\pm$ 5.3 & 0/6 \\
PPO & tuned & raw & nom. & 442k & 4.1 & 150.8 $\pm$ 5.9 & 0/6 \\
\midrule
A2C & tuned & raw & nom. & 367k & 16.5 & 85.2 $\pm$ 28.6 & 0/6 \\
A2C & base & sincos & rob. & 492k & 11.4 & 84.7 $\pm$ 28.4 & 0/6 \\
A2C & base & raw & nom. & 371k & 10.8 & 69.7 $\pm$ 31.3 & 0/6 \\
A2C & tuned & sincos & rob. & 462k & 10.7 & 67.2 $\pm$ 47.6 & 0/6 \\
A2C & tuned & sincos & nom. & 350k & 12.1 & 49.7 $\pm$ 35.6 & 0/6 \\
A2C & base & sincos & nom. & 350k & 13.7 & 39.3 $\pm$ 30.7 & 0/6 \\
\bottomrule
\end{tabular}
\end{table}

\begin{table}[t]
\centering
\caption{Paired comparisons (within-seed blocks). $\Delta$ is the mean paired difference in return; CI is a 95\% percentile bootstrap; $d_z$ is the paired effect size; $p$ is the paired Wilcoxon signed-rank test (\textsuperscript{*}$<$0.05, \textsuperscript{**}$<$0.01, \textsuperscript{***}$<$0.001). Comparisons against a censored side (sin/cos, PPO) are conservative.}
\label{tab:exp1_tests}
\begin{tabular}{l c r c c c}
\toprule
Comparison & $n$ & $\Delta$ & 95\% CI & $d_z$ & Wilcoxon $p$ \\
\midrule
PPO vs A2C (all 36 matched config-seed pairs) & 36 & $+106.6$ & $[+93.4, +119.4]$ & $+2.60$ & $<$0.001\textsuperscript{***} \\
PPO: base reward vs tuned reward & 18 & $+5.0$ & $[+2.0, +8.4]$ & $+0.71$ & 0.027\textsuperscript{*} \\
PPO: sin/cos vs raw obs (nominal only) & 12 & $+23.9$ & $[+19.2, +28.7]$ & $+2.73$ & $<$0.001\textsuperscript{***} \\
PPO: robust vs nominal training (sin/cos only) & 12 & $-1.3$ & $[-3.4, +0.2]$ & $-0.39$ & 1.000 \\
A2C: base reward vs tuned reward & 18 & $-2.8$ & $[-17.8, +12.9]$ & $-0.08$ & 0.859 \\
A2C: sin/cos vs raw obs (nominal only) & 12 & $-32.9$ & $[-54.6, -9.4]$ & $-0.79$ & 0.032\textsuperscript{*} \\
A2C: robust vs nominal training (sin/cos only) & 12 & $+31.5$ & $[+13.3, +50.2]$ & $+0.90$ & 0.084 \\
\bottomrule
\end{tabular}
\end{table}


\paragraph{PPO decisively outperforms A2C.}
PPO reaches a mean return of $172.6$ (median $180.5$) against A2C's $66.0$
(median $92.2$), while training about four times faster ($3.0$ vs.\ $12.6$ min).
The separation is clean: A2C's \emph{best} single run ($137.1$) still falls
below PPO's \emph{worst} ($145.5$), and $15$ of the $36$ A2C runs effectively
fail to learn the swing-up (return $<30$). The advantage is large and
significant in the paired test (\mbox{$\Delta=+106.6$}, 95\%~CI
$[+93,+119]$, $d_z=2.6$, Wilcoxon $p<10^{-10}$; Table~\ref{tab:exp1_tests}), and
since the PPO side is censored at the $180$ cap this gap is in fact a
\emph{lower bound}. A2C is also far less reproducible: a Levene test rejects
equal variance ($p=1.7\times10^{-5}$; std.\ $36.1$ vs.\ $12.3$), with per-config
seed spread reaching a coefficient of variation of ${\sim}78\%$ for A2C against
${\le}4\%$ for PPO. On this task A2C is therefore not competitive.

\paragraph{A censored ceiling near $180$.}
The strong PPO configurations converge to almost exactly the early-stopping cap
($\approx 181$): $22$ of the $72$ runs --- \emph{all} of them PPO\,$\cdot$\,sin/cos
--- hit $180$ and were stopped, so their reported return is right-censored. Final
return therefore cannot rank these configurations against each other (they are
equal by construction), and the only fair discriminator is how \emph{quickly} and
how \emph{reliably} each reaches the cap. On steps-to-threshold the best
controller is \texttt{PPO, tuned reward, sin/cos, nominal}, which reaches $180$ in
$167$k steps ($1.4$ min) on all $6$ seeds.

\paragraph{The sin/cos encoding is the one decisive design factor.}
For PPO, the wrap-invariant encoding yields a real and significant return gain
over the raw angle (\mbox{$\Delta=+23.9$}, 95\%~CI $[+19,+29]$, $d_z=2.7$,
$p<10^{-3}$); because the sin/cos side is censored, the true gain is again at
least this large. Reward shaping is at most second-order for PPO
(\mbox{$\Delta=+5.0$}, $p=0.027$) and irrelevant for A2C. For A2C the sin/cos
encoding actually \emph{hurts} ($\Delta=-32.9$, $p=0.032$), consistent with the
weaker learner struggling with the harder-to-exploit representation.

\paragraph{Robust training is effective but not free.}
Robust-trained PPO still reaches the cap, so it costs no \emph{final return}
(\mbox{$\Delta=-1.3$}, CI $[-3.4,+0.2]$, $p=1.0$) --- but this apparent
``no-cost'' result is an artifact of censoring. On the uncensored metric the
trade is visible: robust training needs more steps to reach $180$ (e.g.\
$167$k\,$\to$\,$342$k for the tuned reward) and, in that setting, only $4/6$
seeds reach the cap at all (vs.\ $6/6$ nominal). Robust training therefore buys
disturbance tolerance for sim-to-real at a sample-efficiency and reliability
cost, rather than for free.

\paragraph{More compute did not help.}
Return correlates negatively with training time ($r=-0.81$ across all runs,
$r=-0.62$ within PPO): the longest runs are the failed A2C runs that trained for
many minutes yet scored low, while the best PPO controllers hit the cap and
stopped within a couple of minutes. Performance is thus set by the algorithm and
the observation encoding rather than by the compute budget, and early stopping
correctly halted converged runs well before the $10^6$-step limit. Reporting
mean~$\pm$~std over six seeds removes the single-seed confound of the earlier
write-up, turning the previously ``confounded'' secondary effects into the
quantified, significance-tested results above.

\begin{figure}[t]
\centering
\includegraphics[width=0.8\linewidth]{figures/algo_comparison.png}
\caption{PPO vs.\ A2C on mean return, training time, and steps trained. Bars
show the per-algorithm mean over all runs and seeds; dots show individual runs.}
\label{fig:exp1_algo}
\end{figure}

\begin{figure}[t]
\centering
\includegraphics[width=0.8\linewidth]{figures/effect_sizes.png}
\caption{Effect sizes of each paired comparison: mean difference in return with
95\% bootstrap confidence intervals. Green = significant gain, red = loss,
grey = interval spans zero (tie).}
\label{fig:exp1_effects}
\end{figure}
